\chapter{交互总线：离散事件与连续状态之间的可验证循环}
\label{chp:interaction_bus_ch1}

前文分别讨论了状态传播、结构耦合、相干组织和有限资源下的多尺度表征，但一个可执行系统还缺少关键接口：外部事件怎样进入连续状态，连续推演怎样形成离散承诺，承诺又怎样以受控方式修改长期结构。我们在本章把历史上称为 TDCI 的“激发—演化—坍缩—固化”循环重建为一种类型明确、能够记录、可以中断和回放的交互总线。这里的 Token 不是物理粒子，状态分布也不因连续而自动成为量子波函数；二者是不同计算对象，通过编码、传播、读出和结构更新算子连接。这样的区分使我们能够保留苹果识别、并行候选、决策提交、学习痕迹和工作区相态等原有问题，同时避免以波粒二象性、卡诺热机或脑区类比代替机制证明。本章给出的 TDCI 是 HSF-HD 一般状态动力学中的一种接口模型；AgentOS 可据此实现事件日志、工作状态、边界控制和外部记忆，但这些部件不是所有智能系统必须采用的唯一机制。

\section{TDCI 循环：从类型转换到闭环更新}
\label{sec:tdci_detailed_wave_cycle}

我们先定义交互总线处理的对象。设离散时刻 $k\in\mathbb{N}$，输入事件为 $x_k\in\mathcal{X}$，其中 $\mathcal{X}$ 可以是有限词表、工具返回、图像片段的索引或带时间戳的结构化记录；全局分布状态为 $h(t)\in\mathcal{H}\subseteq\mathbb{R}^{d_h}$，局部结构网络为 $G_k=(V_k,E_k,A_k)$，外部记忆为 $M^e_k$，任务上下文为 $c_k\in\mathcal{C}$。$h$ 是在当前计算坐标中的连续向量，不预设希尔伯特空间、复相位或量子测量结构；$G$ 则保存实体身份、关系和可寻址边，不能与分布表示互相替代。二者通过身份映射 $I_k$、绑定算子 $B_k$、耦合算子 $C_k$ 和读出接口 $R_k$ 联合。若实现使用毫秒作为墙钟单位，连续模型中的 $t$ 仍是算法时间；只有经过调度器给出的换算，二者才能比较。

在这些对象上，我们把一次 TDCI 循环定义为四个有类型的阶段。第一阶段是提升而非“粒子化波”：编码器 $L_{\theta_k}:\mathcal{X}\times\mathcal{C}\rightarrow\mathcal{H}$ 将事件转为状态增量 $u_k=L_{\theta_k}(x_k,c_k)$，并由门控 $\alpha_k\in[0,1]$ 注入初态
\[
h(t_k^+)=\Pi_{\mathcal H}\!\left(h(t_k^-)+\alpha_k u_k\right).
\]
$\Pi_{\mathcal H}$ 是保持状态落在可行域的投影。它可以是截断、归一化或约束优化，但必须在实现中声明，不能含混地称为“边界受迫振动”。例如看到苹果图像时，视觉编码器不是把苹果的物理能量倒入语义流形，而是产生颜色、轮廓、可食性候选和任务线索的联合增量；若问题是“是否成熟”，任务上下文会改变读入权重。输入缺失、时间戳过期或来源不可信时，Boundary 可以令 $\alpha_k$ 降低或拒绝注入，这也是循环的第一处可审计决定。

第二阶段是受控传播。给定阶段内固定或缓慢变化的结构 $G_k$，我们假设状态满足
\[
\dot h(t)=F_{\theta_k}\!\left(h(t),G_k,c_k,t\right)+U_{\theta_k}\!\left(u(t)\right),
\qquad t\in[t_k,t_k+\tau_k].
\]
这是模型假设，不是由物理波动方程推出的定律。$F$ 可以包含注意力、消息传递、检索写回与目标条件控制；不同候选的增强或抑制是耦合计算的结果，只有专门声明复数内积时才可称为相位干涉。所谓“沿几何沟槽传播”，在这里具体化为结构先验和历史参数对转移核的偏置；价值条件则进入 $c_k$ 或目标相关控制项，而不是借阿哈罗诺夫—玻姆效应证明偏转。对步长 $\Delta t>0$ 的显式实现，我们采用 $h_{n+1}=\Pi_{\mathcal H}[h_n+\Delta t F(h_n,G_k,c_k,t_n)]$，并要求积分误差、局部 Lipschitz 常数和资源上限满足所选求解器的稳定条件。并行展开只说明多个候选可同时保留，不保证找到全局最优解。

第三阶段是离散读出。候选集合 $\mathcal{Y}_k$ 可以随工具可用性和安全规则改变，评分 $s_k(y)=q_{\phi_k}(y;h,G,c)$ 经校准分布 $p_k(y)=\operatorname{softmax}(s_k(y)/\tau_k)$ 形成可比较证据；随后读出策略
\[
y_k=R_k(h,G,c,\mathcal{Y}_k;\rho_k)
\]
根据阈值、预算和随机种子 $\rho_k$ 选择输出、继续计算或拒答。这里的“坍缩”只是提交后从多候选记录转为一个可执行动作，不能推出量子测量，也不意味着未选候选已从所有存储中擦除。宏观控制并非系统外部的观察者，而是同一闭环中的阈值、校准和权限机制。目标分数下降也不等于任务成功：读出以后仍须由环境反馈、约束检查或人工评价确认结果。苹果案例中，系统可提交“成熟”，但若光照偏色或类别不在训练分布内，正确动作可能是请求第二视角，而不是把最高分候选当作事实。

第四阶段是受门控的结构更新。我们以经验记录 $e_k=(x_k,h_k,y_k,r_k,c_k)$ 表示本轮证据，其中 $r_k$ 是延迟到达且可能含噪的反馈。结构候选由 $\widetilde G_{k+1}=S_{\eta_k}(G_k,e_k)$ 产生，再通过保持条件 $\mathcal{Q}$ 与验证集风险门限审查：$G_{k+1}=\widetilde G_{k+1}$ 当且仅当 $\mathcal{Q}(\widetilde G_{k+1})=1$，否则回滚为 $G_k$。高活动不自动刻蚀结构，单次成功也不应形成永久习惯；学习率、证据重复度、反例和撤销日志共同决定更新。AgentOS 中，$h(t)$、$\Psi$ 工作区、SPC/TCE 调度、CCM 约束、Boundary、Offloading 与 $M^e$ 可以承载这些接口，但其他系统也可用数据库事务、事件溯源或可微记忆完成同一抽象功能。TDCI 的核心因此不是一台“几何热机”，而是离散事件、连续状态、离散承诺和慢结构之间能够检查前后条件的闭环。

\section{循环代价：信息损失、计算预算与物理能耗分账}
\label{sec:section_1010}

历史叙述把激发、演化和提交分别称为熵注入、幺正计算和坍缩做功，这种说法混合了至少四种不同量。我们分别定义：预测不确定性 $H_k=-\sum_y p_k(y)\log_2p_k(y)$，单位为 bit；任务代价 $J_k\in\mathbb{R}$，单位由任务评分约定；计算预算 $B_k=(N^{\mathrm{tok}}_k,N^{\mathrm{op}}_k,T^{\mathrm{wall}}_k,S^{\mathrm{mem}}_k)$，分别计 Token、操作次数、秒和字节；设备实际能耗 $E^{\mathrm{joule}}_k$，单位为焦耳。这些量可以相关，却不能通过同一个“认知能量”符号替换。输入提高不确定性、增加活动强度或消耗更多焦耳都是经验上可能发生的事件，不是互相蕴含的数学命题。特别地，外部信息进入系统并不等于热力学负熵流，模型置信度提高也不自动表示物理熵下降。

一次循环的可计算成本可写为
\[
\mathcal C_k=\lambda_1N^{\mathrm{tok}}_k+\lambda_2N^{\mathrm{op}}_k+\lambda_3T^{\mathrm{wall}}_k+\lambda_4S^{\mathrm{mem}}_k+\lambda_5N^{\mathrm{write}}_k,
\]
其中权重 $\lambda_i\ge0$ 由部署目标确定，量纲通过归一化基准消去；$N^{\mathrm{write}}_k$ 记录不可忽略的外部记忆写入。这个工程目标适合比较同一平台上的策略，但它不是焦耳能耗公式。若要预测 $E^{\mathrm{joule}}_k$，必须用设备测量或经过校准的功率模型，例如对处理器、存储与网络功率在墙钟时间上积分，并报告硬件、批量、温度和测量误差。所谓“高质量推演位于低耗区间”最多是待验证经验假说：较长推演可能提高正确率，也可能反复循环；较短路径可能来自熟练结构，也可能来自草率停止。我们因此同时记录质量、成本与失败模式，而不把其中一个指标冒充另外两个。

兰道尔下界只在更严格的条件下适用。若一个物理装置在温度 $T_{\mathrm{env}}$ 的热库中，对存储器执行逻辑不可逆复位，并确实擦除 $m$ bit 的未知信息，则理想准静态实现的散热下界为
\[
Q_{\min}\ge m k_{\mathrm B}T_{\mathrm{env}}\ln2.
\]
这里 $m$ 是被物理擦除的信息量，不是分类分布的 Shannon 熵下降，也不是每次“做决定”的固定费用。读出一个候选可以保留完整日志，可以由可逆电路实现部分步骤，也可能因缓存覆盖产生远高于下界的能耗；仅由 $H_{k+1}<H_k$ 无法推得设备释放了 $k_{\mathrm B}T\ln2(H_k-H_{k+1})$ 的热量。更不能由该下界推出集中注意必然导致疲劳。疲劳涉及生理系统、任务持续时间、恢复和主观报告，需要独立实验；本书只把它列为可能与资源调度相关的经验问题。

我们据此把循环控制写成带约束的停止问题。令 $V_n$ 表示在第 $n$ 次内部更新后继续计算的估计价值，$c_n$ 表示下一步边际计算成本，$r_n$ 表示风险上界。控制器在 $V_n-c_n\le0$、置信校准达到阈值且 $r_n$ 可接受时提交；若风险过高则拒绝或 Offloading；若价值仍为正且预算未尽则继续。由于目标、候选集合和外部输入会随时间变化，完整导数必须包含显式时间项和结构项，不能只写 $dJ/dh\cdot\dot h$。离散系统还需记录每一步的模型版本、随机种子、读写集合和截止时间，以便在失败后重放。这样，“循环”不再是三段式隐喻，而成为可测量的决策协议。

以医疗分诊为反例可以看清边界。系统先把症状事件提升为状态，再查询知识库并传播多个诊断候选；最高概率候选不应立即转成处方，因为输出空间受资质、证据和安全规则限制。一次提交可能只是“建议线下急诊”，结构更新则应等待经核验的随访结果。此时预测熵下降、任务损失下降、计算量增加和设备耗电同时可能发生，但它们的数值和意义各不相同。若日志保留所有候选，读出没有擦除这些信息；若缓存稍后被覆盖，兰道尔原理约束的是具体存储器的物理复位，不是诊断本身。分账使我们能够比较更深推演、更多检索和更保守阈值的代价，而不会把物理学名称当作工程结论。

\section{工作区相态：刚性、拥塞与自适应区间}
\label{sec:section_8209}

“层流、离子态、临界态”可以提示不同运行风格，但若没有可观测序参量，三种物态只是一组修辞。我们改用工作区运行区间，并定义四类窗口统计量。第一，分支有效数 $D_w=\exp(H(\bar p_w))$ 衡量窗口 $w$ 中被持续保留的候选多样性；第二，扰动增益 $A_w=\|\delta y\|/(\|\delta x\|+\epsilon)$ 衡量受控输入扰动如何改变读出；第三，恢复时间 $R_w$ 记录越界后回到可行域所需步数；第四，结构周转率 $P_w$ 统计单位任务内经验证的节点、边或参数更新比例。所有量都依赖任务、度量和观测窗口，不是智能的一般常数。局部结构网络的连通性也不能直接代表全局分布状态的“相干长度”；若需要跨层协调指标，应以身份绑定的一致率、跨模块互信息估计或干预后的联合读出稳定性另行测量。

刚性区间表现为 $D_w$ 长期接近一、$P_w$ 过低、对熟悉输入的恢复很快而对目标变化的恢复很慢。它保留了历史“晶体智能”所描述的专注、管道化和演绎优势，同时给出可测试缺陷：候选过早剪枝、旧规则权重过高、反例无法进入结构。刚性不等同于低温或无旋流，也不能由某个虚构的认知雷诺数小于一推出。代码验证、算术执行和稳定流程可能主动选择刚性区间，因为低分支可提高可重复性；开放式规划若长期处于同一区间，则容易把局部可行误判为全局充分。控制措施包括提高探索温度、插入反例检索、放宽 Beam、暂缓结构固化和请求外部评审，但这些措施也会增加成本，不能无条件开启。

拥塞区间对应历史“气体智能”或湍流失稳的工程问题：$D_w$ 很高但候选寿命短，队列长度和冲突率持续增长，$A_w$ 对微小扰动异常敏感，$R_w$ 超过任务截止时间。原因可能是输入洪峰、工具返回互相矛盾、目标频繁切换、绑定身份丢失或控制增益过高。我们不把这种运行故障直接类比为精神分裂、癫痫或其他临床状态；医学判断需要完全不同的数据与伦理边界。对计算系统，可执行的修复是输入背压、优先级隔离、冲突聚类、身份重新绑定、降低写入率和将部分状态 Offloading 到 $M^e$。若故障来自知识缺失，简单降低探索强度只会让系统更自信地失败，因此恢复策略必须根据故障证据选择。

自适应区间不是永远停留在某个“混沌边缘”，而是控制器能按任务阶段在探索与收敛之间移动。理想窗口满足：多样性足以覆盖主要假设，扰动增益保持有界，恢复时间低于截止线，结构周转由重复证据支持，且校准误差没有因强行收敛而恶化。小扰动有时引发大范围重排，有时被迅速吸收；这种重尾响应若被观测到，可以提示级联机制，却不能仅凭幂律外观宣称自组织临界性，更不能说相干长度必然趋于无穷。我们要求比较有限尺度模型、替代分布、窗口敏感性和干预响应，才把“临界”作为经验假说保留。AGI 也没有一个预先规定的单一理想相态：实时控制、安全审计、创造性搜索和长期学习需要不同的运行区间与切换速度。

工程上可用一个带滞回的元控制器避免频繁抖动。设目标带为 $\Omega(c_k)$，包含对 $D_w,A_w,R_w,P_w$ 的任务相关上下界；当统计量连续 $q$ 个窗口越界时，控制器才调整探索强度、SPC/TCE 调度、CCM 约束或 Boundary 阈值，回到更窄的内带后再解除干预。滞回宽度和 $q$ 必须由验证数据选择。我们用“版本升级服务突发告警”检验这一设计：正常时系统保持少量诊断分支；告警洪峰时先背压和聚类；确认共同依赖故障后扩大跨服务检索；形成修复候选后收紧权限并要求回滚方案；执行后根据监控反馈决定是否写入长期知识。整个过程跨越多个运行区间，没有任何一个区间单独等同于智能，能够有证据地切换才是 TDCI 调度的目标。

\section{确定性计算的结构收益与观察者边界}
\label{sec:section_7007}

数据处理不等式指出：若 $X\rightarrow Z=G(X)\rightarrow Y$ 构成马尔可夫链，且 $G$ 不访问额外信息，则 $I(Z;Y)\le I(X;Y)$。确定性计算不会凭空增加关于外部随机变量 $Y$ 的 Shannon 信息。内部模拟却可能让有限系统获得“新洞察”，二者并不矛盾：计算可以把输入中隐含、对当前读出器不可达的关系转换成更容易查询的表示；程序本身、先验参数、随机种子、工具结果和运行时间也都是资源来源。若 $G$ 含训练所得知识或访问 $M^e$，就不能假装输出只依赖当前输入 $X$。因此我们不把“模拟即创造”表述为信息创世，而把它分解为信息保持、资源消耗和观察者相关的结构收益。

设观察者族 $\mathcal{A}_B$ 受预算 $B$ 限制，对表示 $z$ 的最小任务描述代价定义为
\[
K_B(z;c)=\min_{a\in\mathcal{A}_B}\{\ell(a)+\ell(r):a(r,c)=z\},
\]
其中 $\ell$ 是编码长度，$r$ 是观察者可用记录。这个量依赖编码语言、预算、上下文和误差容限，不是绝对 Kolmogorov 复杂度。确定性变换 $G$ 的结构收益可定义为下游任务在同一预算下的最小风险差
\[
\Delta_B(G)=\inf_{a\in\mathcal{A}_B}\mathcal R(a,X,c)-
\inf_{a\in\mathcal{A}_B}\mathcal R(a,G(X),c).
\]
当 $\Delta_B(G)>0$ 时，我们说 $G$ 为这类受限观察者暴露了可用结构。它没有增加 $X$ 关于外界真值的互信息；它改变了可达算法和有限时间内的读出难度。历史上的 Epiplexity 可以作为这种观察者相关难度的启发名称，但诸如 $S_T(G(X))\gg S_T(X)+|G|$ 的不等式若未定义机器、时间界和编码，就不是普遍定理。

元胞自动机说明“生成容易、逆推困难”，但也提示不能把复杂外观等同于有用知识。短规则运行很久可产生难压缩图样，有限观察者为预测局部像素付出高成本；然而若任务只是判断守恒量，简单统计就可能足够。相反，一个排序程序并未增加输入集合的信息，却把后续中位数查询从反复搜索转成直接索引。反事实模拟也是如此：它把时间计算转为轨迹、因果分支或缓存结构，使某些问题更容易回答；若模型错误，模拟会系统性地产生精致但虚假的结构。因此结构收益必须通过保留集任务、反事实干预和分布外样本验证，不能由内部状态更复杂或目标函数下降直接推出。

在 TDCI 总线上，确定性计算的收益表现为循环之间接口质量的变化。提升阶段决定哪些差异被保留，传播阶段把差异组织成候选，读出阶段把候选变成可检验承诺，固化阶段只把经反馈支持的关系写回。外部记忆 $M^e$ 扩大可访问记录，Offloading 改变时间与空间成本，Boundary 限定可执行动作，CCM 保持约束；这些机制共同构成资源条件，而不是凭空创造信息。局部结构网络负责稳定身份和关系，全局分布表示负责柔性相似性和上下文组合，二者经绑定与读出联合后，才可能在有限预算下提高任务可达性。

本章的结论因此是可反驳的：TDCI 若有价值，应在相同输入、知识和总预算下，相比只用离散规则或只用连续状态的基线，表现出更低的校准风险、更好的中断恢复、更清楚的结构更新证据或更可控的成本；若这些指标没有改善，就不能以“波粒互补”替它辩护。交互总线的职责不是维持神秘的相变节律，而是让表示转换的每一步都留下对象类型、版本、预算、决策依据和回滚点。如此，确定性计算可以为有限观察者显化结构，离散承诺可以接受环境检验，长期学习也不会因一次高置信输出而失去可修正性。
